\documentclass[11pt]{article}

\usepackage[letterpaper,margin=0.8in]{geometry}
\usepackage{amsmath,amssymb}
\usepackage{booktabs}
\usepackage{array}
\usepackage{xcolor}
\usepackage{hyperref}
\usepackage{titlesec}
\usepackage{parskip}
\usepackage{fancyhdr}
\usepackage{enumitem}
\usepackage{multicol}

\hypersetup{
  colorlinks=true,
  linkcolor=blue!50!black,
  urlcolor=blue!50!black,
  pdftitle={LonghornSilicon Compiler Integration Brief},
  pdfauthor={LonghornSilicon},
}

\titleformat{\section}{\normalfont\large\bfseries}{\thesection.}{0.5em}{}
\titleformat{\subsection}{\normalfont\normalsize\bfseries}{\thesubsection}{0.5em}{}

\pagestyle{fancy}
\fancyhf{}
\lhead{\small LonghornSilicon $\times$ Compiler Integration}
\rhead{\small Meeting Brief}
\cfoot{\thepage}
\renewcommand{\headrulewidth}{0.4pt}

\setlist{itemsep=1pt,topsep=2pt,leftmargin=1.2em}

\begin{document}

%==============================================================================
% Header
%==============================================================================
\thispagestyle{empty}
\begin{center}
{\LARGE\bfseries LonghornSilicon $\times$ Silicon-Agnostic Compiler}\\[0.3em]
{\Large\bfseries Integration Brief}\\[0.6em]

\rule{0.8\textwidth}{0.4pt}\\[0.6em]

\begin{tabular}{@{}rl@{\hspace{2em}}rl@{}}
\textbf{Date}:    & 2026-05-13           & \textbf{Branch}:   & \texttt{compiler-integration-prep} \\
\textbf{Contact}: & Chaithu Talasila     & \textbf{Email}:    & \href{mailto:talasilachaithu105@gmail.com}{talasilachaithu105@gmail.com} \\
\textbf{Affil.}:  & UT Austin            & \textbf{Repo}:     & \href{https://github.com/LonghornSilicon/attention-compute-unit}{LonghornSilicon/attention-compute-unit} \\
\end{tabular}\\[0.4em]
\rule{0.8\textwidth}{0.4pt}
\end{center}

\vspace{0.3em}

%==============================================================================
\section{What LonghornSilicon Is Building}
\label{sec:what}

A four-block LLM inference accelerator targeting TSMC 16nm FinFET (N16FFC) via the TSMC
University Program (tape-out target Summer 2027):
\begin{enumerate}
\item \textbf{ACU (Attention Compute Unit)} --- INT8 / FP16 mixed-precision
  attention with a streaming \emph{precision controller} that picks the
  precision per tile from a single-cycle ratio check.\\
  \textbf{Status:} RTL frozen, 253 / 253 tests pass, Sky130 PnR signed off
  (DRC / LVS / antenna / IR-drop clean), ASAP7 area projection, public paper.
\item \textbf{KV Cache Engine} --- on-die SRAM (spilling to off-chip LPDDR5X) with
  ChannelQuant compression (per-channel INT4 keys / per-token INT4 values +
  FP16 outlier lane). \emph{RTL complete through Sky130 sign-off.}
\item \textbf{Token Importance Unit} --- per-token attention-weight
  accumulator driving keep / demote / evict for mixed-precision KV
  retention. \emph{Not yet implemented.}
\item \textbf{Memory Hierarchy Controller} --- routes 0.8 MB on-die SRAM
  $\leftrightarrow$ off-chip LPDDR5X, direct (no separate eDRAM tier).
  \emph{Not yet implemented.}
\end{enumerate}
For this meeting, the precision controller (block 1) is what we can commit
to a stable interface for.

%==============================================================================
\section{What We Offer the Compiler Today}
\label{sec:offer}

Three concrete artifacts the compiler team can target right now, before
silicon or FPGA hardware exist:

\begin{table}[h]
\centering
\small
\begin{tabular}{@{}p{0.32\textwidth} p{0.55\textwidth}@{}}
\toprule
Artifact                                          & Purpose / Status \\
\midrule
Bit-accurate Python reference model\par\hspace*{0.5em}\texttt{sw/reference\_model/precision\_controller\_ref.py}
& A pure-Python class that produces exactly the INT8 / FP16 decision the chip will produce. Verified \textbf{143 / 143 bit-exact} against the RTL testbench. \\[0.4em]
ISA / interface specification\par\hspace*{0.5em}\texttt{docs/isa/precision\_controller\_isa.pdf}
& Versioned (\texttt{pc-isa-0.1}) memory map, AXI protocols, C API stub. Stable for the precision controller; will unify with other blocks as they land. \\[0.4em]
Compiler-perspective demo\par\hspace*{0.5em}\texttt{sw/reference\_model/example\_compiler\_use.py}
& Three usage levels (naive / batched / calibrated) showing what compiler-emitted code looks like against the model. \\
\bottomrule
\end{tabular}
\end{table}

\noindent
A compiler engineer can write a backend against the Python model today
and trust that the chip will agree byte-for-byte when silicon arrives.

%==============================================================================
\section{Four-Phase Integration Plan}
\label{sec:phases}

The interface is stable across all four phases. Same memory map, same
streaming protocol, same model semantics throughout.

\begin{table}[h]
\centering
\small
\renewcommand{\arraystretch}{1.15}
\begin{tabular}{@{}l l p{3.4cm} p{6cm}@{}}
\toprule
Phase                        & Timeline                                      & Compiler targets             & Who delivers what \\
\midrule
0 (Python ref)               & now                                           & Python reference model        & \textbf{Us:} model + ISA + tests already on the branch. \textbf{Compiler team:} backend $\rightarrow$ calls model \\[0.4em]
1 (FPGA prototype)           & when ZCU102 / 104 arrives                     & AXI-Lite + AXI-Stream         & \textbf{Us:} Vivado bitstream + PYNQ driver. \textbf{Compiler team:} backend $\rightarrow$ driver \\[0.4em]
2 (multi-block FPGA)         & after KV cache + token + memory blocks complete & Same AXI, more blocks visible & \textbf{Us:} larger Vivado project. \textbf{Compiler team:} extend backend \\[0.4em]
3 (silicon)                  & 2027+ (post-tape-out)                         & PCIe-attached chip            & \textbf{Us:} chip + Linux driver. \textbf{Compiler team:} re-target backend \\
\bottomrule
\end{tabular}
\end{table}

\noindent
\textbf{Critically:} phase 0 is enough to start. The compiler team does not
need to wait for anything from us beyond what is already on the branch.

%==============================================================================
\clearpage
\section{What We Ask of the Compiler Team}
\label{sec:asks}

\begin{enumerate}
\item A short architecture writeup: what is the IR (MLIR / Relay / custom),
  what is the typical hardware-target shape they integrate against, what
  experience do they have adding new backends.
\item A scoping conversation: is the precision controller the right first
  target, or do they want to wait for more blocks?
\item Honest timeline expectations: phase 0 work could start immediately;
  phase 1 (FPGA) needs the ZCU102 / 104 in hand, which is a separate gating
  item.
\item Their thoughts on the open ISA questions (\S\ref{sec:open}):
  runtime-vs-synthesis-time \texttt{THRESHOLD}, exposing
  \texttt{max} / \texttt{sum} alongside the decision bit, AXI vs other
  fabric.
\end{enumerate}

%==============================================================================
\section{Open Questions for the Meeting}
\label{sec:open}

These are intentionally on the table; we have not committed to them.

\begin{enumerate}
\item Runtime-configurable \texttt{THRESHOLD}? Cost: one register, one
  parameterized multiplier. Benefit: per-workload tuning without
  re-synthesis.
\item Single-bit decision vs.\ exposing \texttt{max} and \texttt{sum}? Lets
  the compiler implement its own gating policy downstream.
\item AXI-Lite + AXI-Stream the right fabric, or CHI / TileLink / custom?
\item Phase 1 timing: do they need real FPGA hardware to begin meaningful
  work, or is the phase 0 Python target enough to start a backend that the
  FPGA later validates?
\end{enumerate}

%==============================================================================
\section{What We Are \emph{Not} Committing To Today}
\label{sec:notcommit}

To keep boundaries clear up front:
\begin{itemize}
\item Schedule promises tied to tape-out --- the TSMC University Program
  shuttle date is outside our control.
\item An ISA for the full chip --- three of the four blocks are not yet
  RTL. The precision-controller ISA is stable; the chip-level ISA is not.
\item Source access to private repos --- until we know the collaboration
  shape, the public repo is the boundary.
\item Exclusivity --- other compiler projects may also be interested in
  this chip. Happy to discuss exclusivity terms, but not by default.
\end{itemize}

%==============================================================================
\section{Concrete Next Step (Roughly One Week)}
\label{sec:next}

If the meeting goes well, the smallest meaningful collaboration looks like:
\begin{enumerate}
\item The compiler team writes a two-page integration plan: which IR ops
  map to which ISA operations, what their backend's runtime expects from
  us, what would unblock them to start coding.
\item We respond within a week with the gaps we can fill and timelines.
\item They pick the first integration target --- likely a single matmul or
  a single attention block exercised against the phase 0 reference model.
\item We co-author the first commit in their backend that produces correct
  decisions when run against our model.
\end{enumerate}

\noindent
That commit becomes the milestone we can both demo at meetings,
conferences, or to UT Austin / TSMC for further support.

%==============================================================================
\vspace{0.5em}
\noindent\rule{\textwidth}{0.4pt}\\
\noindent\small\textit{Everything in this brief, plus the source code,
RTL, OpenLane-signed-off Sky130 GDS, and the precision-controller paper, is
public and shareable on the \texttt{compiler-integration-prep} branch of
\texttt{LonghornSilicon/attention-compute-unit}. The TSMC 16FFC PDK
is NDA-protected and not part of this material.}

\end{document}
